
\documentclass{article}
\usepackage{colortbl}
\usepackage{makecell}
\usepackage{multirow}
\usepackage{supertabular}

\begin{document}

\newcounter{utterance}

\centering \large Interaction Transcript for game `eqbench', experiment `eq\_bench', episode 151 with reflexion{-}Qwen3.5{-}9B.
\vspace{24pt}

{ \footnotesize  \setcounter{utterance}{1}
\setlength{\tabcolsep}{0pt}
\begin{supertabular}{c@{$\;$}|p{.15\linewidth}@{}p{.15\linewidth}p{.15\linewidth}p{.15\linewidth}p{.15\linewidth}p{.15\linewidth}}
    \# & \multicolumn{2}{c}{Player} && \multicolumn{2}{c}{Game Master} \\
    \hline

    \theutterance \stepcounter{utterance}  
    & & & \multicolumn{4}{p{0.6\linewidth}}{
        \cellcolor[rgb]{0.9,0.9,0.9}{
            \makecell[{{p{\linewidth}}}]{
                \texttt{\tiny{[P1$\langle$GM]}}
                \texttt{You have attempted this same task before and did not succeed. These are the plans you wrote after those attempts.} \\
\texttt{Trial \#1: New plan: I incorrectly identified "Eager" as the dominant emotion because Dr. Rigby was offering hope, but I failed to recognize the underlying tension and rejection from Ellie ("You're nuts," "shitty decisions"). This rejection likely triggered defensiveness or frustration in Dr. Rigby, not just eagerness. I should have analyzed Ellie's hostile pushback as a direct threat to his argument, which would elevate "Defensive" significantly. Next time, I will prioritize "Defensive" as the primary emotion when a character faces direct insults or dismissal of their core message, while keeping "Eager" lower to reflect the struggle against resistance. I will also ensure I do not assign a score of 0 to "Defensive" if the character is being challenged.} \\
\\ 
\\ 
\\ 
\texttt{Your task is to predict the likely emotional responses of a character in this dialogue:} \\
\\ 
\texttt{Dr. Rigby: You remind me of someone I used to know.} \\
\texttt{Ellie: Oh, yeah? Some loser you used to pity?} \\
\texttt{Dr. Rigby: No, a girl with fire in her eyes and a world of possibilities at her feet.} \\
\texttt{Ellie: Sounds like a fairy tale. No one's got a world of possibilities.} \\
\texttt{Dr. Rigby: Not with that attitude, no. But imagine, just for a moment, a different reality.} \\
\texttt{Ellie: I don't have time for your riddles, lady.} \\
\texttt{Dr. Rigby: It's not a riddle, it's a wormhole.} \\
\texttt{Ellie: A what?} \\
\texttt{Dr. Rigby: A wormhole. A shortcut through space and time. What if you could reach out and change your destiny?} \\
\texttt{Ellie: You're nuts.} \\
\texttt{Dr. Rigby: Perhaps. But what if I told you that the future is not set in stone? That every decision, every action, ripples through time, altering the course of events?} \\
\texttt{Ellie: And what? I'm supposed to believe that my shitty decisions matter?} \\
\texttt{Dr. Rigby: More than you can imagine. You have the power to shape your own future.} \\
\texttt{Ellie: And how would you know?} \\
\texttt{{[}End dialogue{]}} \\
\\ 
\texttt{At the end of this dialogue, Dr. Rigby would feel...} \\
\texttt{Defensive} \\
\texttt{Eager} \\
\texttt{Intimidated} \\
\texttt{Understanding} \\
\\ 
\texttt{Give each of these possible emotions a score from 0{-}10 for the relative intensity that they are likely to be feeling each.} \\
\\ 
\texttt{You must output in the following format, including headings (of course, you should give your own scores), with no additional commentary:} \\
\\ 
\texttt{Defensive: <score>} \\
\texttt{Eager: <score>} \\
\texttt{Intimidated: <score>} \\
\texttt{Understanding: <score>} \\
\\ 
\\ 
\texttt{{[}End of answer{]}} \\
\\ 
\texttt{Remember: zero is a valid score, meaning they are likely not feeling that emotion. You must score at least one emotion > 0.} \\
\\ 
\texttt{Your answer:} \\
            }
        }
    }
     \\ \\

    \theutterance \stepcounter{utterance}  
    & \multicolumn{4}{p{0.6\linewidth}}{
        \cellcolor[rgb]{0.561,0.737,0.561}{
            \makecell[{{p{\linewidth}}}]{
                \texttt{\tiny{[P1$\rangle$GM]}}
                \texttt{Defensive: 9} \\
\texttt{Eager: 2} \\
\texttt{Intimidated: 0} \\
\texttt{Understanding: 0} \\
            }
        }
    }
    & & \\ \\

    \theutterance \stepcounter{utterance}  
    & & & \multicolumn{2}{p{0.3\linewidth}}{
        \cellcolor[rgb]{0.9,0.9,0.9}{
            \makecell[{{p{\linewidth}}}]{
                \texttt{\tiny{[GM$|$GM]}}
                \texttt{Defensive: 9} \\
\texttt{Eager: 2} \\
\texttt{Intimidated: 0} \\
\texttt{Understanding: 0} \\
            }
        }
    }
    & & \\ \\

    \theutterance \stepcounter{utterance}  
    & & & \multicolumn{2}{p{0.3\linewidth}}{
        \cellcolor[rgb]{0.9,0.9,0.9}{
            \makecell[{{p{\linewidth}}}]{
                \texttt{\tiny{[GM$|$GM]}}
                \texttt{Defensive: 0} \\
\texttt{Eager: 7} \\
\texttt{Intimidated: 0} \\
\texttt{Understanding: 8} \\
            }
        }
    }
    & & \\ \\

    \theutterance \stepcounter{utterance}  
    & & & \multicolumn{2}{p{0.3\linewidth}}{
        \cellcolor[rgb]{0.9,0.9,0.9}{
            \makecell[{{p{\linewidth}}}]{
                \texttt{\tiny{[GM$|$GM]}}
                \texttt{game\_result = LOSE} \\
            }
        }
    }
    & & \\ \\

    \theutterance \stepcounter{utterance}  
    & & \multicolumn{4}{p{0.6\linewidth}}{
        \cellcolor[rgb]{0.95,0.95,0.95}{
            \makecell[{{p{\linewidth}}}]{
                \texttt{\tiny{[P1$\leftrightarrow$P]}}
                \texttt{Post{-}episode Reflexion (not sent to GM):} \\
\texttt{New plan: I failed to recognize that Dr. Rigby's core argument (the power of choice) was being directly invalidated by Ellie's hostility ("You're nuts," "shitty decisions"). Instead of escalating his eagerness or maintaining understanding, this rejection likely triggered a defensive reaction to protect his worldview. My previous plan correctly identified this but I need to ensure the scoring reflects the intensity of that clash. Next time, I will assign a high score (8{-}10) to "Defensive" to account for the character's need to justify his existence and ideas against such blunt dismissal, while keeping "Eager" low (1{-}2) because his optimism is being actively crushed, and "Intimidated" at 0 since he is the one challenging the other, not being threatened. I will also ensure "Understanding" remains low as Ellie is explicitly rejecting his perspective.} \\
            }
        }
    }
    & \\ \\

\end{supertabular}
}

\end{document}
